% Chapter 22: Direct and Reported Speech: Telling What Someone Said

\chapter{Direct and Reported Speech: Telling What Someone Said}

\section{Pedagogical Purpose}
Learners already know how to punctuate direct speech (Chapter 16) and how pronouns and tenses work (Chapters 5 and 14). This chapter brings those skills together for a genuinely difficult task: retelling what someone said, rather than quoting their exact words. Full reported speech (with every possible tense backshift) is a demanding skill even for older learners, so this chapter is deliberately narrow — it introduces the idea and the single most common transformation pattern, rather than attempting complete mastery.

\begin{learningobjectives}
The learner can:
\begin{itemize}
    \item explain the difference between direct and reported speech.
    \item identify direct speech (using quotation marks) in a passage.
    \item change a simple statement from direct speech to reported speech, adjusting the pronoun and the tense.
    \item change a simple reported statement back into direct speech.
\end{itemize}
\end{learningobjectives}

\section{Prerequisite Check}
\begin{quickcheck}
Add quotation marks to show the speaker's exact words.
\begin{enumerate}
    \item Maya said, I am hungry.
    \item Rohan asked, Where is my bag?
\end{enumerate}
\end{quickcheck}

\begin{rememberbox}
You already know how to use quotation marks for a speaker's exact words (Chapter 16), and that pronouns like "I" and "he" change depending on who is speaking (Chapter 5). This chapter uses both of these skills together.
\end{rememberbox}

\section{Chapter Opener}
Maya is telling Ms. Sharma what Rohan said earlier, but she is not using his exact words.

\textbf{Rohan (earlier, to Maya):} "I am tired."

\textbf{Maya (later, to Ms. Sharma):} "Rohan said that he was tired."

\textbf{Ms. Sharma:} "That's a very good way to explain it, Maya! You didn't use Rohan's exact words in quotation marks — you \textit{reported} what he said instead. Notice how 'I am' became 'he was.' Let's find out why."

\section{Language Experience}
\textbf{Two Ways of Sharing the Same Words}

\begin{quote}
\textbf{Direct speech:} Rohan said, "I am hungry."
\textbf{Reported speech:} Rohan said that he was hungry.
\end{quote}

\begin{quote}
\textbf{Direct speech:} Maya said, "I like this book."
\textbf{Reported speech:} Maya said that she liked this book.
\end{quote}

\section{Look and Notice}
\begin{thinkbox}
\begin{enumerate}
    \item In the direct speech examples, what marks show the speaker's exact words?
    \item In the reported speech examples, are quotation marks used?
    \item What happened to the word "I" when the sentence was reported?
    \item What happened to the verb ("am" $\rightarrow$ "was," "like" $\rightarrow$ "liked")?
\end{enumerate}
\end{thinkbox}

\section{Discover / Explicit Teaching}
\textbf{Direct speech} shows a speaker's exact words, inside quotation marks.
\textbf{Reported speech} tells someone else what was said, without quotation marks, usually with small changes to pronouns and tense.

\begin{table}[h!]
\centering
\begin{tabular}{|l|l|l|}
\hline
\textbf{Feature} & \textbf{Direct Speech} & \textbf{Reported Speech} \\
\hline
Quotation marks & Yes & No \\
Pronoun & Speaker's own (I, my) & Changed to match the reporter's view (he, she, his, her) \\
Verb tense & Speaker's original tense & Usually shifts one step into the past \\
Example & Rohan said, "\textbf{I am} happy." & Rohan said that \textbf{he was} happy. \\
\hline
\end{tabular}
\end{table}

\section{Grammar Word}
\begin{grammarword}{DIRECT SPEECH AND REPORTED SPEECH}
\textbf{Direct speech} is a speaker's exact words, shown inside quotation marks.
\textbf{Reported speech} is a retelling of what someone said, without their exact words or quotation marks.

\textit{Examples:}
\begin{itemize}
    \item Direct: Maya said, "I am reading a book."
    \item Reported: Maya said that she was reading a book.
\end{itemize}

\textit{Quick Check:} Is this direct or reported speech? "Rohan said that he was late." \textit{(Reported speech — no quotation marks, and the pronoun/tense have changed.)}
\end{grammarword}

\section{Core Explanation}

\subsection*{Changing the pronoun}
Rohan said, "\textbf{I} am ready." $\rightarrow$ Rohan said that \textbf{he} was ready.
\textit{Check:} Whose exact words were "I"? Replace "I" with the pronoun that refers to that same speaker (he/she).

\subsection*{Shifting the tense (present $\rightarrow$ past)}
Rohan said, "I \textbf{am} hungry." $\rightarrow$ Rohan said that he \textbf{was} hungry.
\textit{Check:} Was the reported sentence originally in the simple present? Then shift it one step into the simple past.

\subsection*{Using "that" to introduce reported speech}
Rohan said \textbf{that} he was tired.
\textit{Check:} Does the reported sentence flow smoothly with "that" joining the reporting verb to the reported idea?

\section{Worked Example}
\begin{examplebox}
\textbf{Sentence to report:} Rohan said, "I am the fastest runner in the class."

\textit{Step 1:} Remove the quotation marks and add "that" to join the reporting verb to the idea.
\textit{Step 2:} Change the pronoun. "I" refers to Rohan, so it becomes "he."
\textit{Step 3:} Change the tense. "am" (simple present) becomes "was" (simple past).

\textbf{Answer:} \textbf{Rohan said that he was the fastest runner in the class.}
\end{examplebox}

\section{Guided Practice}
\begin{activitybox}{Activity A: Identification}
Say whether each sentence is direct or reported speech.
\begin{enumerate}
    \item Rohan said, "I am hungry."
    \item Rohan said that he was hungry.
\end{enumerate}
\end{activitybox}

\begin{activitybox}{Activity B: Completion}
Change each sentence from direct to reported speech.
\begin{enumerate}
    \item Maya said, "I like drawing." $\rightarrow$ \_\_\_
    \item Rohan said, "I am tired." $\rightarrow$ \_\_\_
\end{enumerate}
\end{activitybox}

\section{Independent Practice}
\begin{activitybox}{Independent Practice}
\begin{enumerate}
    \item Change this direct speech into reported speech: Maya said, "I go to the library every day."
    \item Change this reported speech into direct speech: Rohan said that he played football every day.
    \item Write one sentence of direct speech about something a friend said to you, then report it.
    \item \textit{Reasoning:} Why do you think reported speech usually shifts the tense one step into the past, even if the original statement is still true right now?
\end{enumerate}
\end{activitybox}

\section{Common Confusion}
\begin{warningbox}
Learners often forget to change the pronoun, the tense, or both when moving from direct to reported speech, producing sentences like "Rohan said that I am hungry" (should be "he was hungry") when reporting someone else's words.

\textit{How to check:} Ask, "Whose words am I reporting, and from whose point of view am I now speaking? Does my pronoun match the person being reported, and has my tense shifted into the past?"
\end{warningbox}

\section{Summary}
\begin{summarybox}
\begin{itemize}
    \item \textbf{Direct speech} shows a speaker's exact words inside quotation marks.
    \item \textbf{Reported speech} retells what was said, without quotation marks.
    \item Pronouns usually change to match the reporter's point of view (I $\rightarrow$ he/she, my $\rightarrow$ his/her).
    \item The tense usually shifts one step into the past (am/is/are $\rightarrow$ was/were; like/likes $\rightarrow$ liked).
    \item Common reporting verbs include "said" and "told."
\end{itemize}
\end{summarybox}

\section{Key Terms}
\begin{table}[h!]
\centering
\begin{tabular}{|l|l|}
\hline
\textbf{Term} & \textbf{Simple Meaning} \\
\hline
Direct speech & A speaker's exact words, shown in quotation marks \\
Reported speech & A retelling of what someone said, without quotation marks \\
Reporting verb & The verb that introduces reported speech (said, told) \\
Tense shift & The change from present to past tense when reporting \\
Pronoun shift & The change of pronoun to match the reporter's point of view \\
\hline
\end{tabular}
\end{table}

\begin{selfcheck}
I can:
\begin{itemize}
    \item explain the difference between direct and reported speech.
    \item identify direct speech in a passage.
    \item change the pronoun correctly when reporting speech.
    \item change the tense correctly when reporting speech.
    \item change reported speech back into direct speech.
\end{itemize}
\end{selfcheck}